\section{結言}
この日のXでは、強い専門能力を持つモデルの公開と、その能力を実務に組み込むインターフェースの整備が並行して進んだ。GPT-6 Cyberの発表は安全保障上の利用範囲と高度な研究能力の両立を問い、ML Internはモデル開発の手順をチャットから操作する方向を示した。

研究投稿では、モデルの真正性を検証する暗号技術や、シミュレーションと実データの整合、分子・CAD設計の評価基盤が取り上げられた。モデルの性能そのものに加えて、検証可能性、再現性、実行環境の構成が重要になっている。

一方、次期モデルの公開予測や個人による速度・ベンチマーク比較には未確認の主張が含まれる。観測窓外の投稿も一部収録したため、速報性と確度の境界を意識しながら読む必要がある。今後は限定公開モデルの利用範囲拡大と、公開された研究・ツールの実装詳細が継続的な注目点となる。
